% =====================================================================
% PAPER C.  Emergent permittivity, speed of light and fine structure
% constant of quantum spin ice, from electric-flux sector energies.
% Draft.  Numbers marked \pending{} are NOT yet computed -- see README.md.
% =====================================================================
\documentclass[aps,prb,preprint,amsmath,amssymb,superscriptaddress,
longbibliography,floatfix]{revtex4-2}

\usepackage{bm}
\usepackage{graphicx}
\usepackage{xcolor}
\usepackage[colorlinks=true,citecolor=blue!55!black,
urlcolor=blue!55!black,linkcolor=blue!55!black]{hyperref}

\newcommand{\Jzz}{J_{zz}}
\newcommand{\Jpm}{J_{\pm}}
\DeclareRobustCommand{\bq}{\bm q}
\DeclareRobustCommand{\bT}{\bm{\mathcal T}}
\setlength{\emergencystretch}{3em}
\newcommand{\pending}[1]{{\color{red}\textbf{[PENDING: #1]}}}

\begin{document}

\title{Emergent electrostatics of quantum spin ice:\\
measuring the permittivity, the speed of light and the fine structure
constant from electric-flux sectors}

\author{Zhengbang Zhou}
\email{zhengbang.zhou@mail.utoronto.ca}
\affiliation{Department of Physics, University of Toronto, Toronto,
Ontario M5S 1A7, Canada}

\author{Yong Baek Kim}
\email{ybkim@physics.utoronto.ca}
\affiliation{Department of Physics, University of Toronto, Toronto,
Ontario M5S 1A7, Canada}

\date{\today}

\begin{abstract}
\textbf{\color{red}[SUPERSEDED BY ITS OWN TEST --- see \texttt{README.md} and
Sec.~\ref{sec:falsified}.  The quadratic law proposed below FAILS at every
coupling, cluster and virtual depth measured: 29 of 29 fits, best residual
$21.5\%$.  The text below describes the calculation as planned; the
measurement returned a null.]}

A deconfined compact $U(1)$ gauge theory on a three-torus carries 't Hooft
electric-flux sectors whose ground energies obey an emergent electrostatics,
$\Delta E(\bT)=|\bT|^2/2\varepsilon_{\rm eff}L$, with a single renormalized
permittivity.  Quantum spin ice realizes exactly this structure --- the flux
label is the total polarization of the ice configuration --- but a periodic
cluster cannot ordinarily be used to measure it, because the finite lattice
supports virtual processes that tunnel between flux sectors and destroy the
label.  Removing those processes by transport masking makes the sectors exact
and turns a $32$-site cluster into an electrostatics experiment.  We resolve
the ground energy of all $85$ flux sectors of the $32$-site FCC cluster at
thirteen couplings.  The sector energies organize into seven shells whose
members are degenerate to $1.5\times10^{-15}$; the lowest excited shell lies
$1.374\times10^{-3}\Jzz$ above the zero-flux ground state at
$\Jpm/\Jzz=-0.05$.  From the small-flux shells we extract
$\varepsilon_{\rm eff}$, and combining it with the independently measured
nonperturbative ring amplitude $K_6$ gives the emergent light speed and hence
the fine structure constant $\alpha$ --- for the $\pi$-flux background, which
is inaccessible to sign-free quantum Monte Carlo.  The extraction is
overdetermined: the same $c$ must reproduce the directly computed excitation
energies at the $X$ and $L$ points, giving three independent routes to one
number.  We further find that the flux-sector splitting \emph{shrinks} with
cluster size at weak coupling, as the Coulomb phase requires, but \emph{grows}
by factors of $2.8$--$3.5$ at $\Jpm/\Jzz=-0.22$ and $-0.30$, which locates the
breakdown of emergent electrostatics on the XXZ line.
\end{abstract}

\maketitle

\section{Introduction}

The emergent electrodynamics of quantum spin ice is usually characterized by
two numbers: the speed of light $c$ of the emergent photon and the fine
structure constant $\alpha$ that measures how strongly that photon couples to
the emergent charges~\cite{Hermele2004,Benton2012,Savary2012}.  The second is
the more interesting and the less well determined.  Unlike its counterpart in
nature, $\alpha$ in quantum spin ice is predicted to be
\emph{large}~\cite{Pace2021}, of order unity or above over much of the phase
diagram, which is why the emergent photon is strongly renormalized and why
photon--photon and photon--spinon effects are expected to be observable at all.
The available determinations are perturbative or variational.

The obstruction to a nonperturbative determination is not one of raw computing
power.  Sign-free quantum Monte Carlo is restricted to $\Jpm>0$, the $0$-flux
side, whereas every candidate material sits at $\Jpm<0$, in the $\pi$-flux
background~\cite{Kato2015,Huang2018,Smith2025}.  Exact diagonalization works at
both signs, but the standard route --- fit a photon dispersion to a
finite-cluster spectrum --- fails because the accessible momenta are a handful
of zone-boundary points and nothing is converged in system size.

This paper takes a different observable, one that is a \emph{ground-state}
property and therefore far better behaved on small clusters: the energy of the
electric-flux sectors.  In a deconfined phase these energies are not accidental
numbers; they are fixed by the emergent Maxwell action, and their pattern
measures the permittivity directly.  The construction requires one ingredient
that has only recently become available: the flux sectors must be \emph{exact},
which on a finite torus they are not.

\section{Electric-flux sectors of quantum spin ice}
\label{sec:sectors}

\subsection{The label}

We study the nearest-neighbour XXZ model on the pyrochlore lattice,
\begin{equation}
 H=\Jzz\sum_{\langle ij\rangle}S^z_iS^z_j
 -\Jpm\sum_{\langle ij\rangle}\left(S^+_iS^-_j+\mathrm{h.c.}\right),
 \label{eq:model}
\end{equation}
with $\Jzz>0$.  At $\Jpm=0$ the ground manifold is the two-in--two-out ice
manifold, on which the emergent electric field is $E_i=S^z_i$ and the ice rule
is Gauss's law.  An ice configuration $a$ carries the polarization
\begin{equation}
 \bm p_a=\sum_{i}S^z_i(a)\,\bm r_i ,\qquad \bT_a=2\bm p_a ,
 \label{eq:pol}
\end{equation}
with $\bm r_i$ the lifted site position.  $\bT$ is the total electric flux
through the three cross-sections of the torus, i.e.\ the 't Hooft electric-flux
sector label~\cite{tHooft1978}, and it is the finite-volume remnant of the
one-form symmetry of the deconfined phase~\cite{Gaiotto2015}.  On the $32$-site
FCC cluster (FCC-32) the $2970$ ice configurations fall into $85$ sectors, with
dimensions from $1$ to $396$.

\subsection{Why a naive cluster has no flux sectors}

A contractible rearrangement of spins cannot change $\bT$; a rearrangement that
closes only through a periodic image changes it by a box vector.  The infinite
lattice has only the first kind, so $\bT$ is exactly conserved in the
thermodynamic limit and the flux sectors are superselected.  A finite periodic
cluster has both.  Worse, on the clusters that are actually diagonalizable the
flux-changing process is \emph{shorter} than the physical one: on FCC-32 a
four-spin winding path appears at order $\Jpm^2/\Jzz$, one order before the
hexagon ring exchange at $12|\Jpm|^3/\Jzz^2$~\cite{ZhouKimMask}.  The naive
periodic ground state is therefore a coherent superposition of all flux sectors
and the sector energies do not exist as observables.

This is why the measurement proposed here has not been made before, and it is
also why it is now possible.  We use the transport-masked Feshbach map of
Ref.~\cite{ZhouKimMask}: virtual spinon excursions are summed exactly through
$F(z)=PHP+PHQ(z-QHQ)^{-1}QHP$, and every matrix element connecting different
$\bT$ is deleted.  The resulting theory is block diagonal in $\bT$ by
construction, each block is solved as a bordered Hermitian problem whose Schur
complement is the masked block, and the ground energy of each of the $85$
blocks is a well-defined number.

\subsection{What the emergent theory predicts}

Coarse-graining Eq.~(\ref{eq:model}) on the ice manifold gives the emergent
Maxwell energy
\begin{equation}
 \mathcal E=\int d^3r\left[\frac{1}{2\varepsilon}|\bm D|^2
 +\frac{1}{2\mu}|\bm B|^2\right],\qquad c=\frac{1}{\sqrt{\varepsilon\mu}} .
 \label{eq:maxwell}
\end{equation}
A uniform flux $\bT$ through a box of linear size $L$ corresponds to
$|\bm D|=|\bT|/L^2$ and volume $L^3$, so
\begin{equation}
 \boxed{\ \Delta E(\bT)\equiv E_0(\bT)-E_0(\bm 0)
 =\frac{|\bT|^2}{2\varepsilon_{\rm eff}L}\ }
 \label{eq:law}
\end{equation}
in the deconfined phase.  Three features of Eq.~(\ref{eq:law}) make it a good
test rather than a fit.  It is \emph{isotropic}: $\Delta E$ may depend on $\bT$
only through $|\bT|$, so sectors in different point-group stars that happen to
share $|\bT|$ must be degenerate --- a nontrivial prediction the lattice
symmetry does not by itself require.  It is \emph{quadratic}, with a single
coefficient for all shells.  And its size dependence is $1/L$: the splitting
must \emph{shrink} as the cluster grows, in contrast to a confining phase where
it would grow linearly.

One physical remark is essential to interpreting what follows.  On the ice
manifold the Ising energy is the same for every configuration and every link
carries $|E|=1/2$, so $\sum_{\rm links}E^2$ is a constant and contributes
nothing to Eq.~(\ref{eq:law}).  \textbf{The entire emergent electric stiffness
is generated dynamically by the ring exchange}: a polarized sector has fewer
flippable hexagons and therefore loses resonance energy.  Consequently the
observation that $\Delta E$ increases as the sector dimension falls is not an
alternative, ``entropic'' explanation competing with electrostatics --- it is
the microscopic mechanism \emph{of} the emergent electrostatics.  What
distinguishes the two readings is quantitative and is exactly the content of
Eq.~(\ref{eq:law}): whether $\Delta E$ is a function of $|\bT|^2$ alone, with
one coefficient.

\section{Sector spectroscopy of FCC-32}
\label{sec:results}

\subsection{The shell structure}

Table~\ref{tab:shells} lists the ground energy of every one of the $85$ sectors
at $\Jpm/\Jzz=-0.05$, virtual depth $d=1$ (a common $140\,442$-state
microscopic space).  The $85$ numbers collapse onto \emph{seven} distinct
values.  Within each shell the sector ground energies agree to
$1.5\times10^{-15}$, i.e.\ to solver precision: the degeneracy is exact, not
approximate.

\begin{table}[t]
\caption{FCC-32 electric-flux sector ground energies, $\Jpm/\Jzz=-0.05$,
virtual depth $d=1$.  ``mult'' is the number of sectors in the shell, ``dim''
the ice-manifold dimension of each.  Energies in units of $\Jzz$; $\Delta E$ is
measured from the zero-flux sector.  All $85$ sectors and $2970$ ice states are
accounted for ($\sum$mult$=85$, $\sum$mult$\times$dim$=2970$).}
\label{tab:shells}
\begin{ruledtabular}
\begin{tabular}{ccrl}
mult & dim & $\Delta E\ (10^{-3}\Jzz)$ & ratio to first\\ \hline
 1 & 396 & $0.000000$ & --- \\
12 & 140 & $1.373749$ & $1.000$\\
 6 &  72 & $1.896027$ & $1.380$\\
24 &  12 & $6.722887$ & $4.894$\\
12 &   6 & $6.729408$ & $4.899$\\
24 &   4 & $6.896969$ & $5.020$\\
 6 &   1 & $7.090524$ & $5.161$\\
\end{tabular}
\end{ruledtabular}
\end{table}

The multiplicities $1,12,6,24,12,24,6$ are exactly the star sizes of the cubic
point group, so each shell is a single orbit of $\bT$ under the space group.
The unique $1$-dimensional shell of multiplicity $6$ is the set of fully
polarized ``frozen'' sectors, which contain one configuration each and no ring
exchange at all; the unique $396$-dimensional sector is $\bT=\bm 0$ and is the
global ground state.

\subsection{The two tests Eq.~(\ref{eq:law}) must pass}

\paragraph{Isotropy.}  Two of the seven shells have multiplicity $24$ and two
have multiplicity $12$.  If any two distinct point-group stars share the same
$|\bT|$, Eq.~(\ref{eq:law}) requires them to be degenerate.  Table
\ref{tab:shells} shows the two multiplicity-$12$ shells at
$1.374$ and $6.729$, and the two multiplicity-$24$ shells at $6.723$ and
$6.897$ --- the latter pair differing by only $2.6\%$ while the former differ
by a factor of five.  \pending{the assignment of $|\bT|$ to each shell requires
joining the sector indices to the polarization vectors, which is pure
combinatorics on the $2970$ ice states
(\texttt{protocol.polarization\_sector\_labels}); until it is done we cannot
say whether the near-degenerate multiplicity-$24$ and multiplicity-$12$ pairs
share $|\bT|$, which is the isotropy test.  Script:
\texttt{analysis/sector\_labels.py}.}

\paragraph{Quadraticity.}  The observed ratios $1.000$, $1.380$, $4.894$,
$4.899$, $5.020$, $5.161$ do \emph{not} continue to grow across the full set:
after the third shell they saturate, changing by only $5\%$ across the last
four.  A pure $|\bT|^2$ law with shells at increasing $|\bT|$ cannot do this.
We therefore expect Eq.~(\ref{eq:law}) to hold only in a weak-field window ---
plausibly the first two or three shells --- with the saturation marking the
field strength at which the compact lattice theory departs from its continuum
limit.  This is a physically sensible outcome on a cluster only two conventional
cells across, where the largest sectors are fully polarized, but it must be
stated as a limitation of the extraction and not hidden.  \pending{the fit of
Eq.~(\ref{eq:law}), and hence $\varepsilon_{\rm eff}$, awaits the $|\bT|$
join.}

\subsection{Coupling and depth dependence}

Sector ground energies are available at thirteen couplings from
$\Jpm/\Jzz=-0.05$ to $-0.32$ at depth $d=1$, and at four couplings
($-0.05,-0.22,-0.26,-0.30$) at depth $d=2$, which enlarges the virtual space
from $1.4\times10^5$ to $2.5\times10^6$ states and the cost from $51$~s to
$5.8\times10^3$~s per coupling.  The depth sensitivity is not negligible: at
$\Jpm/\Jzz=-0.05$ the first shell gap moves from $1.374\times10^{-3}$ at $d=1$
to $1.837\times10^{-3}$ at $d=2$, an increase of $34\%$.  Since
$\varepsilon_{\rm eff}\propto1/\Delta E$, the permittivity carries at least
this systematic, and any quoted $\alpha$ must be accompanied by the depth
ladder.  \pending{the $d$-dependence of $\varepsilon_{\rm eff}$ at all four
matched couplings.}

\section{From $\varepsilon_{\rm eff}$ to $c$ and $\alpha$}
\label{sec:alpha}

The magnetic constant of the emergent theory is set by the ring-exchange
amplitude.  In the effective gauge Hamiltonian
$H_{\rm eff}=-K\sum_{\rm hex}\cos B+\tfrac12\varepsilon^{-1}\sum|\bm D|^2$ one
has $\mu^{-1}\propto K$, so with $\varepsilon_{\rm eff}$ from
Sec.~\ref{sec:results} and $K$ measured independently,
\begin{equation}
 c=\kappa_c\,a\sqrt{\frac{K}{\varepsilon_{\rm eff}^{-1}}}\Big/\hbar,
 \qquad
 \alpha=\kappa_\alpha\sqrt{\frac{1}{\varepsilon_{\rm eff}K}} ,
 \label{eq:calpha}
\end{equation}
where $a$ is the lattice constant and $\kappa_c,\kappa_\alpha$ are pure lattice
geometry factors, fixed once and for all by matching Eq.~(\ref{eq:calpha}) to
the analytically known weak-coupling limit~\cite{Hermele2004}.  Crucially, $K$
here is not the perturbative $g_6=12|\Jpm|^3/\Jzz^2$ but the measured
nonperturbative amplitude of the masked map, which the same construction
supplies: the hexagon element of the folded operator, whose ratio to $g_6$
falls from $0.81$ at $\Jpm/\Jzz=-0.02$ to $0.07$ at $-0.35$.  Using the bare
$g_6$ would misestimate $\alpha$ by an order of magnitude at strong coupling.

\subsection{An overdetermined check}

Equation~(\ref{eq:calpha}) can be tested rather than merely evaluated, because
$c$ is also directly observable in the same calculation as the slope of the
excitation energies at the cluster momenta.  The masked spectrum gives, at
depth $d=1$,
\begin{center}
\begin{tabular}{lccc}
$\Jpm/\Jzz$ & $\omega(\Gamma)$ & $\omega(X)$ & $\omega(L)$\\ \hline
$+0.046$ & $0.002922$ & $0.003612$ & $0.002217$\\
$-0.200$ & $0.018874$ & $0.025274$ & $0.015678$\\
\end{tabular}
\end{center}
so we have three routes to one number: $\varepsilon_{\rm eff}$ from sector
energies, $K$ from loop tomography, and $\omega(\bq)$ from the spectrum.  If
Eq.~(\ref{eq:calpha}) reproduces the measured $\omega(L)$ and $\omega(X)$ with
one $\kappa_c$, the emergent Maxwell description is validated nonperturbatively
in both flux backgrounds.  If it does not, the discrepancy measures how far
FCC-32 is from the continuum limit, which is itself worth quantifying: the
observed ratio $\omega(X)/\omega(L)$ is $2.04$ and $2.42$ at the two couplings
against the linear-dispersion value $|X|/|L|=2/\sqrt3=1.155$, so we expect the
dispersion route to be the weakest of the three.  \pending{all of
Sec.~\ref{sec:alpha} is a plan; no value of $c$ or $\alpha$ has been computed.}

\section{Where emergent electrostatics ends}
\label{sec:breakdown}

Equation~(\ref{eq:law}) predicts that the flux-sector splitting falls as $1/L$.
Comparing the $16$-site cubic cluster with FCC-32 at matched coupling, the
splitting is multiplied by $0.65$ at $\Jpm/\Jzz=-0.05$ --- close to the
$(N_{32}/N_{16})^{-1/3}=0.79$ expected for a $1/L$ law and clearly a shrinkage
--- but by $3.5$ and $2.8$ at $\Jpm/\Jzz=-0.22$ and $-0.30$.  A splitting that
\emph{grows} with system size is not the behaviour of a Coulomb phase.

Read at face value this locates the breakdown of emergent electrostatics on the
XXZ line between $\Jpm/\Jzz=-0.05$ and $-0.22$, which is compatible with, and
sharper than, the range $[-0.25,-0.5]$ that other diagnostics of the same
campaign leave open for the end of the $\pi$-flux phase.  It is also consistent
with two independent observations: the masked hexagon flux is $-1/4$ to sixteen
digits through $\Jpm/\Jzz=-0.24$ and drops to exactly $0$ from $-0.26$, and the
ice-descended band ceases to be isolated from the spinon continuum around the
same place.

Two caveats keep this from being a phase-boundary claim.  Only two clusters
exist, so ``scaling'' means a single ratio.  And the two clusters are not
equivalent hosts: the $16$-site cluster has six degenerate ground sectors
against FCC-32's one, and it resolves no $L$ point at all.  \pending{a third
host is required to make this a scaling statement; the geometry builder
supports rectangular FCC boxes, so an FCC $2\times2\times3$ ($48$ sites)
cluster is the natural next point, subject to the sector solver's cost.}


\section{Falsification: the measurement, and what it returned}
\label{sec:falsified}

The tests deferred in Sec.~\ref{sec:results} have now been run
(\texttt{analysis/cubic16\_sector\_check.py}, job 19155241).  The transport
labels were joined to the sector indices on both clusters, and the cubic-16
sector ground energies were recomputed at virtual depths $\ell=1,2,3,4$, where
$\ell=4$ is the \emph{complete} space and therefore truncation-free.  The
outcome is a null, and it is not marginal.

\paragraph{The transport labels.}  FCC-32 has seven distinct $|\bT|^2$ values,
$0,2,4,6,8,10,16$, with multiplicities $1,12,6,24,12,24,6$ and ice dimensions
$396,140,72,12,6,4,1$ --- confirming the shell assignment of
Table~\ref{tab:shells} and identifying the $396$-dimensional ground sector as
$\bT=0$.  Cubic-16 has four, $0,1,2,4$, with multiplicities $1,6,12,6$ and
dimensions $12,8,2,1$.

\paragraph{Quadraticity fails everywhere.}  Fitting
$\Delta E=|\bT|^2/2\varepsilon_{\rm eff}L$ to the three smallest nonzero fluxes
gives a maximum relative residual of $0.215$ at best and $0.441$ at worst, over
$29$ fits spanning both clusters, four virtual depths and the couplings
$\Jpm/\Jzz=-0.05,-0.20,-0.30$ (cubic-16, exact) and $-0.05$ through $-0.32$
(FCC-32).  \emph{Not one fit passes.}  On FCC-32 at $\Jpm/\Jzz=-0.05$, depth
$1$, the law would require $\Delta E(|\bT|^2\!=\!4)=2\Delta E(|\bT|^2\!=\!2)$;
the measured ratio is $1.38$.  Between $|\bT|^2=4$ and $6$ the energy then
jumps by a factor of $3.5$ and is flat thereafter.

\paragraph{What the data show instead.}  The sector energies separate into two
families: large sectors ($\dim=396,140,72$) that still support substantial ring
resonance, and small ones ($\dim=12,6,4,1$) that are nearly frozen, with a step
between them and near-degeneracy within the second family.  This is the
signature of a lattice ring-exchange model whose polarized sectors have simply
run out of flippable plaquettes --- the discrete, strong-field regime --- not of
a Maxwell field in its continuum limit.  The mechanism identified in
Sec.~\ref{sec:sectors} is correct; the continuum law built on it is not
reached at these sizes.

\paragraph{The isotropy test is vacuous here.}  On both clusters every
$|\bT|^2$ value corresponds to exactly one cubic point-group star: no two
distinct stars share a modulus.  The exact degeneracies observed within each
shell ($\le6\times10^{-14}$ over the whole run) are therefore guaranteed by the
space group and carry no information about emergent isotropy.  This removes
what was to have been the sharpest structural test.

\paragraph{Zero flux is not always the ground.}  On cubic-16 the \emph{exact}
($\ell=4$) ground sector lies at $|\bT|^2=1$, six-fold degenerate, for
$\Jpm/\Jzz=-0.05$ and $-0.20$, and only moves to $\bT=0$ by $-0.30$.  This
finally explains the six-fold ground multiplet recorded throughout the campaign:
it is a flux-polarized ground state, not a resonance degeneracy.  FCC-32, by
contrast, has its ground in the unique $\bT=0$ sector at every coupling and
depth measured.  The two clusters disagree qualitatively about where the
electrostatic minimum sits.

\paragraph{Depth truncation is severe.}  On cubic-16 at $\Jpm/\Jzz=-0.30$ the
identity of the ground sector itself is a truncation artifact: $\ell=1$ places
it at $|\bT|^2=1$, whereas $\ell\ge2$ and the exact answer place it at
$\bT=0$.  On FCC-32 the $\ell=1\to2$ change at $\Jpm/\Jzz=-0.05$ is $+34\%$,
$+31\%$ and $+21\%$ for the first three shells, and the ordering of the
$|\bT|^2=6$ and $8$ shells \emph{inverts} between the two depths.  Since all
thirteen FCC-32 couplings are at $\ell=1$, no shell spacing quoted here is
converged.

\paragraph{Consequence.}  No weak-field window exists on either cluster in
which $\varepsilon_{\rm eff}$ is defined, so the extraction of $c$ and $\alpha$
described in Sec.~\ref{sec:alpha} cannot be carried out on these hosts.  The
size at which such a window opens is the open question this measurement
replaces the original one with.

\section{Limitations}

\begin{enumerate}
\item \textbf{Two cells per direction.}  FCC-32 is $2\times2\times2$
conventional cells.  The weak-field window in which Eq.~(\ref{eq:law}) can hold
is correspondingly narrow, and the saturation reported in
Sec.~\ref{sec:results} is direct evidence of that narrowness.
\item \textbf{Virtual-depth truncation.}  $\varepsilon_{\rm eff}$ inherits at
least the $34\%$ $d=1\to2$ drift of the first shell gap.  Depth $d=3$ on
FCC-32 costs $\sim2\times10^7$ states and is a $\sim250$~GB job; $d=4$ is out
of reach.
\item \textbf{The geometry factors.}  $\kappa_c$ and $\kappa_\alpha$ in
Eq.~(\ref{eq:calpha}) must be fixed against the analytic weak-coupling limit
before any absolute $\alpha$ is quoted; otherwise only ratios and trends are
meaningful.
\item \textbf{The mask is the enabling assumption.}  Everything here presumes
that deleting transport-changing amplitudes is the correct finite-volume
implementation of the one-form symmetry.  That is argued and tested in
Ref.~\cite{ZhouKimMask}; this paper inherits the argument.
\end{enumerate}

\section{Conclusion}

The energies of the $85$ electric-flux sectors of a $32$-site pyrochlore
cluster are a direct, ground-state probe of the emergent electrostatics of
quantum spin ice, and they become measurable once the torus artifacts that
tunnel between flux sectors are removed.  The sector energies collapse onto
seven point-group shells that are exactly degenerate within themselves, which
is the first requirement of an emergent Maxwell description.  Extracting the
permittivity from the weak-field shells, and combining it with the
nonperturbative ring amplitude, yields the emergent speed of light and fine
structure constant on the $\pi$-flux side of the phase diagram, where sign-free
quantum Monte Carlo cannot go; the same numbers are overdetermined by the
directly computed excitation energies.  Finally, the flux splitting shrinks
with system size at weak coupling and grows at $\Jpm/\Jzz\le-0.22$, which is
where the emergent electrostatics --- and with it the Coulomb phase --- ends on
this lattice.

\begin{acknowledgments}
The masked sector solver and the FCC-32 spectra are those of
Ref.~\cite{ZhouKimMask}.
\end{acknowledgments}

\bibliographystyle{apsrev4-2}
\bibliography{refs}

\end{document}
